\documentclass[12pt]{article}
\usepackage[utf8]{inputenc}
\usepackage[T1]{fontenc}
\usepackage{amsmath, amssymb}
\usepackage{geometry}
\usepackage{hyperref}
\usepackage{graphicx}
\usepackage{lmodern}
\usepackage{microtype}
\usepackage{newunicodechar}

% Map Unicode typography to LaTeX-safe ASCII equivalents for stable PDF builds
\newunicodechar{’}{'}
\newunicodechar{“}{``}
\newunicodechar{”}{''}
\newunicodechar{–}{--}
\newunicodechar{—}{---}

\geometry{margin=1in}
\emergencystretch=2em
\hbadness=10000
\hypersetup{
  colorlinks = true,
  linkcolor = blue,
  citecolor = blue,
  urlcolor = blue
}

\begin{document}

\title{Rank from References: \\Operational Evaluation for Text AI Classification}
\author{Zeno Loi, Mircea Sofonea, Grégoire Mercier, Kévin Yauy, David Morquin\\
CHU Montpellier, Univ. Montpellier}
\date{2026-04-21}

\maketitle

\begin{abstract}
\textbf{Background:}\\
Text-classification evaluations often rely on a single reference label, yet in clinical text annotation, competent experts frequently disagree for legitimate reasons. With large language models (LLMs) making operational text structuring plausible, there is a need for \emph{variability-aware} acceptance tests that answer: \emph{Is the AI comparable to its human expert peers?}

\textbf{Methods:}\\
We introduce \textit{Rank from References} (RFR), a parameter-light, auditable acceptance rule. RFR compares an automated classifier to each expert via Hamming distances over binary multi-label outputs and aggregates these comparisons with a rank statistic and document-level bootstrap confidence intervals to implement a quantile-based two-CI non-inferiority/superiority ladder. We additionally report normalized micro-RFR as a consensus percentile for continuous interpretation. We pair RFR with a clinician-grounded simulator that models base prevalence, text-driven co-activation, institutional calibration, annotator-specific thresholds, and within-shift drifts (linear trends and session-reset random walks).

\textbf{Results:}\\
Across \(>2.1\)M simulated classifications, RFR and F1-based methods were used to distinguish unbiased from biased automated classifiers. RFR prioritized safety (higher specificity/precision) over exhaustivity (recall) compared with F1-based rules, with comparable accuracy and modest Matthews correlation coefficient (MCC), reflecting task difficulty. RFR reached specificity \(45.9\%\), precision \(58.2\%\), recall \(60.7\%\), and F1 \(54.9\%\) (vs.\ Average pairwise/Majority F1 specificity \(21.9\%/20.5\%\), precision \(55.3\%/53.7\%\), recall \(85.0\%/81.8\%\), F1 \(63.4\%/61.3\%\)). Error-profile analysis showed substantially fewer false acceptances of biased systems under RFR. Additional sensitivity analyses on panel size (\(m\)) and non-inferiority quantile (\(q\)) showed that \(q\) is the dominant permissiveness lever, while \(m\) effects are context-dependent and should be checked explicitly.

\textbf{Interpretation:}\\
RFR provides a transparent acceptance criterion aligned with safety-first deployment: accept an AI only when its disagreement pattern remains within empirically observed expert variability at a clinically pre-specified human quantile. Consensus-percentile reporting complements this binary decision with an interpretable continuous score. The approach complements training-time robustness methods, supports ablations, and opens the path to evaluating richer text structuring targets (e.g., FHIR-like forms) by substituting Hamming with semantics-aware distances. Given the stakes, variability-aware acceptance testing of this kind should be integrated into pre-deployment evaluation of AI text classifiers that can affect patient care.
\end{abstract}


\newpage
\tableofcontents
\newpage

\section{Quick Start \& Formulas}
\subsection{General Notation}
\begin{itemize}
    \item $\mathcal{D} = \{d_1, \dots, d_{n_d}\}$: A set of $n_d=|\mathcal{D}|$ documents.
    \item $\mathcal{C} = \{c_1, \dots, c_{n_c}\}$: A set of $n_c=|\mathcal{C}|$ categories.
    \item $\Omega = \{0,1\}^{n_d \times n_c}$: The space of \textbf{all} possible classifications.
\end{itemize}

A \textbf{classification} $x \in \Omega$ is represented by a binary vector $(x_{j,k})$, where $x_{j,k} = 1$ if document $d_j$ is assigned to category $c_k$, and $x_{j,k} = 0$ otherwise.

\noindent $\chi = \{h_1, \dots, h_m\} \subset \Omega$ denotes a \textit{collective} set of $m$ human experts, considered as a \textbf{non-unique} gold standard.

\noindent We use the \textbf{Manhattan (Hamming) distance}, denoted $\mathrm{d}(x,y)$. For $x,y \in \Omega$:
\begin{equation}
\mathrm{d}(x,y) = \sum_{j=1}^{n_d} \sum_{k=1}^{n_c} \bigl|x_{j,k} - y_{j,k}\bigr|.
\end{equation}

\subsection{Rank from References}

\subsubsection{Metric Definition: $\rho_{\chi}(x)$}
For each annotator $h_i \in \chi$, we determine the rank of $\mathrm{d}(h_i,x)$ among $\{\mathrm{d}(h_i,h_j)\}_{j=1, j\neq i}^m \cup \{\mathrm{d}(h_i,x)\}$. Denote this value $\rho_i(x)$. We then aggregate:
\begin{equation}
\rho_{\chi}(x) = \frac{1}m\sum_{i=1}^m \rho_i(x).
\end{equation}
\noindent \textbf{Interpretation:} $\rho_{\chi}(x)$ is the \textit{mean rank} of $\mathrm{d}(h_i,x)$ relative to each annotator $h_i$’s distances to other annotators. A higher $\rho_{\chi}(x)$ indicates that, for many annotators, $x$ lies among the greatest observed divergences in $\chi$.

Rank from References can be computed directly from pairwise Hamming distances among all participants (the humans and the candidate classifier). Subsequent sections illustrate how these values inform a go/no-go decision in sensitive domains like healthcare.

\subsubsection{Mean Disagreement Rate Estimator: Distance from References $\Delta_{\chi}(x)$}

This additional metric represents the absolute distance from references to complement the interpretation of $\rho_{\chi}(x)$.
\begin{equation}
\Delta_{\chi}(x) = \frac{1}{m}\sum_{i=1}^{m} \frac{\mathrm{d}(x, h_i)}{\,n_d \times n_c\,}.
\end{equation}
\noindent \textbf{Interpretation:} This is the \textit{mean distance} (weighted by $1/m$) between $x$ and the $m$ annotators, normalized to $[0,1]$ by dividing by $n_d \times n_c$. It indicates the mean disagreement rate between the evaluated classifier and human experts.

\section{Background}
\subsection{Context and Literature Review}
Algorithmic performance in text classification typically assumes a \textit{gold standard} is both known and correct. Yet, in many domains---especially classification tasks demanding subjective human judgment (e.g., patient feedback post-hospitalization)---the uniqueness of any single gold standard is suspect \cite{dawid1979maximum, hui1980study}. Divergences among annotators can be random or reflect distinct interpretive perspectives. Correspondingly, researchers have explored:
\begin{enumerate}
\item \textbf{Latent-label inference} (Dawid \& Skene, 1979; Hui \& Walter, 1980).
\item \textbf{Bayesian aggregation} (Chou \& Lee, 2019; Fornaciari et al., 2021), modeling annotator distributions.
\item \textbf{Robust methods} for noisy labels (co-teaching, loss correction, etc.; Garg et al., 2021).
\item \textbf{Embedding-based nearest-neighbor} techniques (Di Salvo et al., 2024).
\end{enumerate}
Additionally, Fr\'enay \& Verleysen (2014) studied the impact of \textit{label noise} on learning; Cannings et al. (2018) examined theoretical bounds under labeling uncertainty; Bo Han et al. (2022, 2023) investigated large-scale imperfect data. Basile et al. (2021) advanced \textbf{perspectivist data}, endorsing multiple "valid truths." In \textit{medical} settings, Shojania \& Marang-van de Mheen (2020) emphasized that a rigid gold standard for patient safety can be partly arbitrary or incomplete, given legitimate disagreements between experts.


\subsection{The Need for New Metrics}
Despite significant progress in label modeling and error correction, we lack \textbf{practical metrics} for deciding if an AI classifier, in the absence of a singular consensus, is close enough to \textit{human-level} variability. In high-stakes scenarios like healthcare, one needs a \textbf{clear} threshold to judge acceptance:
\begin{quote}
\emph{"Does this AI’s classification pattern remain within the acceptable range demonstrated by our recognized experts?"}
\end{quote}
As Shojania \& Marang-van de Mheen (2020) argue, choosing whether to trust or reject an algorithmic system must be transparent and justifiable.
In parallel, hospital adoption of AI—including LLM-based text tools—has accelerated, yet prevailing governance documents emphasize lifecycle management and general safety principles rather than \emph{variability-aware acceptance testing} for subjective, multi-annotator text tasks. Recent guidance and evaluations (WHO on LMMs, FDA draft guidance for AI-enabled software, NHS lessons from real-world AI evaluations, and international consensus on trustworthy AI) underscore the need for transparent, auditable evaluation but stop short of specifying metrics that test whether an AI’s disagreement pattern stays within expert variability envelopes \cite{who2025lmm,fda2025draft,nhs2024eval,futureai2025}. At the same time, the disruptive potential of LLMs makes operational text structuring newly plausible in clinical workflows, amplifying the risks of deploying systems judged only by consensus-centric metrics such as F1 \cite{tam2024humanEval,singhal2023clinKnowledge,reddy2024genAI}. We address this gap because, in subjective clinical text settings, reliance on traditional point metrics alone remains insufficient for deployment decisions; our work instead advances an acceptance criterion explicitly tied to observed human–human variability.



\subsection{Required Characteristics}
\begin{enumerate}
\item \textbf{Respect for multiple annotators} (no forced consensus).
\item \textbf{Interpretability} for diverse stakeholders (practitioners, regulators, engineers).
\item \textbf{A decisive criterion} (go/no-go) for real-world acceptance.
\item \textbf{Minimized hyperparameters}, avoiding arbitrary thresholds.
\item \textbf{Binary-data suitability}, clarifying results in purely discrete classification tasks.
\end{enumerate}

\subsection{Toward Novel Metrics}
In response, we propose two complementary metrics---\(\rho_{\chi}(x)\) and \(\Delta_{\chi}(x)\)---rooted in a \textbf{collective} gold standard \(\chi\). Section~\ref{sec:methods} details these metrics and shows how they align with the demands of real-world, high-stakes classification.

\subsection{Scope of Application}
\label{sec:scope}

This work targets the evaluation of automated classifiers in multi-annotator clinical text settings under safety-first constraints. Precisely, the intended scope is:

\paragraph{Task type.}
Binary (boolean) \emph{per-category} document classification with \emph{non-exclusive} labels (multi-label). Each category is a yes/no decision at the document level, and an annotator may assign a document to any subset of categories (including none or all).

\paragraph{Input modality.}
Unstructured clinical narratives (e.g., resuscitation reports, progress notes, discharge summaries). Structured fields may be present in the record but are not required and are treated only insofar as they appear in text.

\paragraph{Output format.}
For each document $i$ and category $k$, the system produces either (i) a hard label $\hat{y}_{i,k}\in\{0,1\}$ or (ii) a calibrated probability $\hat{p}_{i,k}\in[0,1]$ subsequently thresholded to a hard label. Ranking-only outputs without a decision threshold are out of scope.

\paragraph{Label schema.}
Categories are \emph{not} mutually exclusive and may co-activate. No hierarchical or ordinal constraints are assumed. If such constraints exist, they are not enforced by default in the present framework.

\paragraph{Unit of analysis.}
Document-level decisions; sequence tagging (token/span-level), multi-instance (bag-of-notes), and section-specific classification are out of scope for the core methodology.

\paragraph{Annotation setting.}
Multiple human experts provide labels with possible institutional differences and individual drifts over time. The framework explicitly models and leverages this variability; single-annotator gold standards are not assumed.

\paragraph{Evaluation target.}
Assessment focuses on whether model disagreement with humans stays \emph{within} the empirically observed envelope of human–human variability. Acceptance is framed via a safety-oriented rule (non-inferiority to expert variability), not solely by point metrics.

\paragraph{Data regime.}
Applicable from small-to-moderate corpora (tens to low thousands of documents) up to large simulated grids. The approach does not require category balance and tolerates heterogeneous prevalences.

\paragraph{Model class.}
Our work is compatible with any text classifier producing per-category decisions or probabilities, including rule-based systems, classical ML, and LLM-based pipelines. No dependence on a specific architecture.

\paragraph{Deployment context.}
We aim clinical decision-support and quality-monitoring scenarios where false acceptance of biased behavior carries asymmetric risk. Cross-site evaluation is supported through explicit modeling of institutional calibration.

\paragraph{Out of scope.}
Exclusive multiclass classification, regression/score prediction, causal inference, active learning policies, and real-time interactive labeling are not addressed here. Extensions to hierarchical labels, multilingual corpora, span extraction, and temporal aggregation are possible but not covered in this study.


\section{RFR Epistemological Construction}\label{sec:methods}
\subsection{Axiomatic Foundation and Spaces}
\begin{enumerate}
\item \textbf{Baseline space} $\Omega$: \\
   Let $\mathcal{D}$ be a set of $n_d$ documents and $\mathcal{C}$ a set of $n_c$ categories. Then $\Omega = \{0,1\}^{n_d \times n_c}$ represents \textbf{all possible} binary classification assignments.
\item \textbf{Collective gold standard} $\chi$: \\
   Instead of identifying a single ground truth, define $\chi = \{h_1,\dots,h_m\}\subset\Omega$ as an \textit{ensemble} of human experts, each considered credible. \newline \textbf{Axiom}: every $h_i\in\chi$ is sufficiently reliable such that pairwise disagreements among them are normal, reflecting the natural spread of opinions rather than outlier behavior.
\item \textbf{Hamming distance}: \\
   For $x,y\in\Omega$,\\
   \begin{equation}
   \mathrm{d}(x,y) = \sum_{j=1}^{n_d} \sum_{k=1}^{n_c} |x_{j,k}-y_{j,k}|.
   \end{equation}
   A straightforward count of how many bits (document--category associations) differ between $x$ and $y$. \newline \textit{Rationale}: Hamming distance in $\{0,1\}$ is highly interpretable, aligns with purely discrete categories, and avoids complexities inherent to continuous embeddings or weighting.
\end{enumerate}

\subsection{W: Human-Expert Space Construction and Approximation}\label{sec:Wconstruction}

The set \( W \) represents the \textit{theoretical space of acceptable classifications}, capturing the inherent variability and pluralism of expert annotations. Formally, it is constructed as a subset of the total classification space \( \Omega = \{0,1\}^{n_d \times n_c} \), constrained by the empirical range of human disagreement. Unlike conventional gold-standard approaches that assume a single correct classification, we define \( W \) such that any classification \( x \in \Omega \) belongs to \( W \) if and only if, for every expert \( h_1 \in \chi \), $x$ is no farther from $h_1$ than at least one other expert \( h_2 \in \chi \).

\begin{equation}
W = \Bigl\{ x \in \Omega \mid \forall h_1 \in \chi, \; \exists h_2 \in \chi, h_2 \neq h_1 , \mathrm{d}(h_1, x) \leq \mathrm{d}(h_1, h_2) \Bigr\}.
\end{equation}

This formulation ensures that \( W \) is \textit{empirically grounded}, encompassing all classifications that lie within the natural scope of human expert disagreement. Theoretically, \( W \) is a \textit{non-parametric, data-driven set} in high-dimensional binary classification space, where boundaries are defined by observed human variability rather than arbitrary thresholds. Classifications within \( W \) are therefore \textit{indistinguishable from human-level variation}, while classifications outside \( W \) exhibit divergences beyond known inter-annotator disagreement. Thus, \( W \) provides a \textit{pragmatic and operational basis for acceptable classification uncertainty}, balancing mathematical rigor with practical interpretability in high-stakes decision-making domains. In essence, \emph{"a classification $x$ is acceptable if, for every expert $h_1\in\chi$, $x$ is closer to $h_1$ than at least one other expert is"}.

\begin{enumerate}
\item \textbf{Alternatives}: spheres or convex hull\\
   One might define $\mathrm{d}(h,x)\le\theta$ or a convex hull of $\chi$. Both typically require an extra parameter $\theta$ and risk over-/under-inclusion. Our approach draws solely on \textbf{actual distances} among the $\chi$ annotators, introducing \textbf{no} hyperparameter.
\item \textbf{Interpretation}:\\
   Provided $\chi$ contains no extreme outlier, the maximum divergences among its members bound $\mathrm{d}(h_1,h_2)$. If $x$ never exceeds that bound for any expert $h_1$, it sits \textbf{within} the domain of plausible human variation.
\end{enumerate}



\subsection{Panel-relative Metrics: Distance and Rank}

\paragraph{Setup.}
Let $d\in\{1,\dots,n_d\}$ index documents, $c\in\{1,\dots,n_c\}$ index categories, and $i,j\in\{1,\dots,m\}$ index experts.
Denote by $y^{(i)}_{d,c}\in\{0,1\}$ the label of expert $i$ for document $d$ and category $c$, and by $x_{d,c}\in\{0,1\}$ the system's label.
Let $\mathrm{d}_d(u,v)=\sum_{c=1}^{n_c}\lvert u_{d,c}-v_{d,c}\rvert$ be the per-document Hamming distance and $\delta_d(u,v)=\mathrm{d}_d(u,v)/n_c$ its normalized version.

\subsubsection{Distance from Reference: $\Delta_{\chi}(x)$}
\begin{enumerate}
\item \textbf{Per-document mean distance:}
\begin{equation}
\Delta_d(x)\;=\;\frac{1}{m}\sum_{i=1}^{m}\delta_d\!\big(y^{(i)},x\big)
\;=\;\frac{1}{m\,n_c}\sum_{i=1}^{m}\mathrm{d}_d\!\big(y^{(i)},x\big).
\label{eq:per-doc-delta-fused}
\end{equation}

\item \textbf{Corpus-level (global) distance:}
\begin{equation}
\Delta_{\chi}(x)\;=\;\frac{1}{n_d}\sum_{d=1}^{n_d}\Delta_d(x)
\;=\;\frac{1}{m\,n_d\,n_c}\sum_{d=1}^{n_d}\sum_{i=1}^{m}\mathrm{d}_d\!\big(y^{(i)},x\big).
\label{eq:global-delta-fused}
\end{equation}

\item \textbf{Justification \& interpretation:}
All experts are equally weighted. The mean (rather than median) preserves localized extremes.
Normalization by $n_c$ (per document) and averaging over $n_d$ makes $\Delta_{\chi}(x)\in[0,1]$ and comparable across corpora:
$\Delta_{\chi}(x)\approx 0$ indicates near-perfect alignment; $\Delta_{\chi}(x)\approx 1$ indicates near-orthogonal labeling.
\end{enumerate}

\subsubsection{Rank from Reference: $\rho_{\chi}(x)$}
\begin{enumerate}
\item \textbf{Per-document, per-expert rank:}
For expert $i$ and document $d$, define the comparison multiset
\[
\mathcal{S}_{i,d}\;=\;\big\{\mathrm{d}_d\!\big(y^{(i)},y^{(j)}\big):\, j\neq i\big\}\ \cup\ \big\{\mathrm{d}_d\!\big(y^{(i)},x\big)\big\}.
\]
Let $\mathrm{rank}_{\mathrm{avg}}(t;\mathcal{S})$ be the ascending \emph{average} rank of $t$ within $\mathcal{S}$ (ties get the average of occupied ranks).
Then
\begin{equation}
r_{i,d}(x)\;=\;\mathrm{rank}_{\mathrm{avg}}\!\Big(\mathrm{d}_d\!\big(y^{(i)},x\big)\,;\,\mathcal{S}_{i,d}\Big)\ \in\ \{1,\dots,m\}.
\label{eq:per-doc-per-exp-rank-fused}
\end{equation}

\item \textbf{Per-document panel-average rank and optional normalization:}
\begin{equation}
\rho_d(x)\;=\;\frac{1}{m}\sum_{i=1}^{m} r_{i,d}(x),
\qquad
\tilde{\rho}_d(x)\;=\;\frac{\rho_d(x)-1}{m-1}\in[0,1].
\label{eq:per-doc-rho-fused}
\end{equation}

\item \textbf{Corpus-level (global) rank:}
\begin{equation}
\rho_{\chi}(x)\;=\;\frac{1}{n_d}\sum_{d=1}^{n_d}\rho_d(x),
\qquad
\tilde{\rho}_{\chi}(x)\;=\;\frac{\rho_{\chi}(x)-1}{m-1}\in[0,1].
\label{eq:global-rho-fused}
\end{equation}

\item \textbf{Display convention \& interpretation:}
We report
\[
\rho_{\chi}(x)\,/\,m\ \ [\min_{i,d} r_{i,d}(x)\ ;\ \max_{i,d} r_{i,d}(x)]
\]
for readability (or the normalized $\tilde{\rho}_{\chi}(x)$). If $\max_{i,d} r_{i,d}(x)\le m-1$, the system is \emph{never} farther from an expert than that expert's farthest peer on any document.
\end{enumerate}

\paragraph{Link to the acceptable set $W$.}
Let $\chi=\{h_1,\dots,h_m\}\subset\Omega$ and let $\mathrm{d}$ be the Hamming distance on $\Omega$.
Define
\[
W \;=\; \Bigl\{x\in\Omega:\ \forall i,\ \exists j\neq i,\ \mathrm{d}(h_i,x)\le \mathrm{d}(h_i,h_j)\Bigr\}.
\]
For each $i$, define the (corpus-level) per-reference rank
\[
\rho_i(x)\;=\;1+\#\bigl\{j:\ \mathrm{d}(h_i,h_j)<\mathrm{d}(h_i,x)\bigr\}\ \in\ \{1,\dots,m\}.
\]
Then:
\begin{enumerate}
\item If $x\in W$, we have $\max_i \rho_i(x)\le m-1$.
\item Conversely, if $\exists i, \rho_i(x) = m$, then $x\notin W$.
\end{enumerate}
The proof follows by noting that $j\neq i$ is excluded from self-comparisons and that at least one peer must be as far or farther than $x$ for membership in $W$.

\subsubsection{Confidence Intervals (Document-Level Aggregation)}
We assume between-document variability dominates within-document category variability.
Accordingly, for a summary metric $\theta\in\{\Delta_{\chi},\rho_{\chi}\}$, define the per-document values
$\theta_d$ (i.e., $\theta_d\in\{\Delta_d,\rho_d\}$) and their mean $\bar{\theta}=\tfrac{1}{n_d}\sum_{d=1}^{n_d}\theta_d$.
A nonparametric bootstrap over documents is then used to quantify uncertainty: we resample
the $n_d$ documents with replacement, recompute $\bar{\theta}$ for each bootstrap replicate
$b\in\{1,\dots,B\}$, and obtain $\bar{\theta}^{*(b)}$.
Hence a $(1-\alpha)$ two-sided confidence interval is estimated with percentile bounds:
\[
\boxed{\;\Bigl[q_{\alpha/2}\bigl(\{\bar{\theta}^{*(b)}\}_{b=1}^{B}\bigr),\ 
q_{1-\alpha/2}\bigl(\{\bar{\theta}^{*(b)}\}_{b=1}^{B}\bigr)\Bigr]\;}\tag{CI}
\]
where, for a $95\%$ interval, the bounds are the $2.5^{\text{th}}$ and $97.5^{\text{th}}$
percentiles of the bootstrap distribution.
This bootstrap formulation is simple to audit and focuses uncertainty on case-to-case
variability across documents without relying on normality assumptions.







\subsection{Practical Verification of the Main Hypothesis}
\begin{enumerate}
\item \textbf{Assumption}: $\chi$ excludes any complete outlier.
\item \textbf{Validation}:\\
   One might check inter-annotator agreement (e.g., Cohen’s kappa, ICC). If real maximum distances among $\chi$ are extremely large, $W$ may become broad, so ensuring expert quality in $\chi$ is crucial.
\end{enumerate}
In summary:
\begin{itemize}
\item \textbf{No model-tuned} acceptance hyperparameters.
\item \textbf{One clinical calibration parameter} ($q$), pre-specified from deployment constraints.
\item \textbf{Reflects} genuine variation in $\chi$.
\item \textbf{Derives} two metrics---$\Delta_{\chi}$ and $\rho_{\chi}$---positioning a new classifier within the inter-annotator space.
\end{itemize}


\subsection{Decision Rule: Non-Inferiority and Superiority}
\begin{enumerate}
    
\item \textbf{Step 1 --- Quantile-based non-inferiority via two confidence bounds.}

    Let $(L_{AI}, U_{AI})$ be the normalized Rank-from-References (lower is better) \((1-\alpha)\) bootstrap CI.

    For each human expert $h_i$, compute the upper \((1-\alpha)\) bootstrap CI bound \(U_i\) by treating $h_i$ as an external candidate against $\chi\setminus\{h_i\}$. Let $\{U_i\}_{i=1}^m$ denote these human reference values, and define the non-inferiority threshold as the clinically pre-specified quantile:
\[
T_q = Q_q\big(\{U_i\}_{i=1}^m\big), \qquad q\in(0,1].
\]
    Non-inferiority holds if:
\[
L_{AI} < T_q\quad\text{and}\quad
U_{AI} < T_q. \tag{Two one-sided tests at level $\alpha$}\footnote{We avoid a Kolmogorov--Smirnov two-sample test because tied discrete $\rho_{\chi}$ ranks violate its continuity assumptions and it does not provide a directly interpretable tolerance margin in this context \cite{conover1972}.}
\]

\item \textbf{Step 2 --- Superiority (optional).} \\

    If non-inferiority is satisfied, we test superiority against a concrete quantile anchor expert using a paired bootstrap on per-document rank differences. Let \(i_q=\arg\min_i |U_i-T_q|\) (human closest to the quantile threshold), with deterministic tie-break by smallest expert index. The anchor \(i_q\) is fixed from the original sample (not re-selected inside bootstrap resamples), and \(\rho_d^{(q)}\) denotes this anchor human's document-level rank trajectory when evaluated externally against \(\chi\setminus\{h_{i_q}\}\).

\[
\delta_d=\rho_d^{AI}-\rho_d^{(q)}, \qquad \bar{\delta}=\frac{1}{n_d}\sum_{d=1}^{n_d}\delta_d.
\]

Let \([L_\delta,U_\delta]\) be the percentile \((1-\alpha)\) CI of \(\bar{\delta}\) from document-level bootstrap resampling.
Superiority is declared if \(U_\delta<0\).

    Algorithmic \textbf{superiority} is declared only when both non-inferiority and this bootstrap condition hold.

\item \textbf{Practical phrasing for non-statisticians.}\\
    The model is accepted if its disagreement profile stays below a clinically chosen human reference quantile; it is called \emph{better than humans} when it also outperforms this quantile anchor in paired document-level comparison.
\end{enumerate}

\section{Evaluative Simulation}
To evaluate the ability of $\rho_{\chi}$ to discriminate intra-$W$ AIs from biased ones, we conducted an in~silico evaluation against widely used alternatives: Majority-F1 and Average pairwise F1.

\subsection{Model Choice}

\paragraph{Qualitative grounding from clinician interviews.}
Following a pilot labeling session on $n=50$ resuscitation reports across $k=11$ categories with $m=6$ intensivists, unstructured interviews surfaced five recurrent themes:
(i) \emph{uneven base rates} (“some flags are almost never present, others are ubiquitous”);
(ii) \emph{text-driven cues and co-occurrence patterns} (“certain wordings almost force you to tick two boxes together”);
(iii) \emph{institutional calibration} (“our unit is stricter on \textit{X} than others”);
(iv) \emph{individual thresholds that drift within a shift} (“after twenty charts my bar changes—fatigue, time pressure, or warming-up”); and
(v) \emph{annotator-specific noisiness} (“some colleagues are consistently sharper on \textit{A} than \textit{B}”).
Clinicians also emphasized that each category is a binary decision at the point of annotation (checked vs.\ not checked), even when categories represent graded clinical phenomena. These observations called for a binary, multi-label simulator that is interpretable in terms of odds and thresholds, yet flexible enough to capture temporal drift, institutional effects, and latent text structure.

\paragraph{Why a logistic, additively-separable linear predictor.}
A logit link offers a clinically intuitive \emph{odds-ratio} interpretation and supports additive contributions on the log-odds scale. Additivity mirrors how clinicians described their reasoning: a baseline tendency modulated by text evidence, then nudged up or down by institutional habits and personal thresholds. We therefore use a logistic model, with each term directly motivated by interview themes:

\begin{itemize}
\item \textbf{Base prevalence}. Captures category-specific baseline odds noted by clinicians (theme i), ensuring realistic marginal prevalences.
\item \textbf{Text effect.} Latent text factors plus residual encode “wording that forces boxes together” (theme ii), allowing shared linguistic cues and category co-activation without hand-crafted rules.
\item \textbf{Institutional-group effect.} Doctors reported systematic unit-level calibration differences (theme iii); fixed contrasts per institution reproduce these shifts in thresholds across services.
\item \textbf{Annotator$\times$category drift.} Fixed bias captures stable personal thresholds (theme v); a linear trend reflects warming-up or fatigue across a session (theme iv); a session-reset random walk adds non-stationary fluctuations clinicians described (“bar changes”).
\item \textbf{Annotator global drift.} Some tendencies are cross-category, being globally stricter/lenient as the shift progresses (themes iv–v).
\item \textbf{Annotator–category scale.} Interviews highlighted that some annotators are not only biased but also \emph{noisier} on specific categories. A scale factor permits heteroscedastic variability and calibration differences beyond mean shifts; the deterministic setting sets $s_{a,k}=1$, while stochastic regimes can inflate or shrink variability as observed in practice.
\end{itemize}

\paragraph{Justification of each hyperparameter family.}
The hyperparameter set is included for use—not fixed a priori—because each governs a phenomenon repeatedly mentioned by clinicians. Their \emph{values} are subsequently estimated from data; their \emph{presence} is dictated by the qualitative evidence.

\paragraph{Alternatives considered and reasons for rejection.}
We examined several modeling families:

\begin{itemize}
\item \emph{Static Dawid–Skene / item-difficulty only.} These EM-based annotator models capture confusability but assume stationarity; they cannot reproduce within-shift drift (themes iv–v) without substantial extensions, and they do not natively encode text-driven co-activations (theme ii).
\item \emph{IRT (Rasch/2PL/3PL) and ordinal IRT.} Powerful for \emph{single} latent trait per item; our multi-label setting with category-specific institutional and temporal effects requires many traits and time dynamics, undermining interpretability and identifiability at $n=50$.
\item \emph{Probit GLMM.} Statistically similar to logit. In our setting, the logistic link simplified calibration to observed odds and yielded more stable simulation control.
\item \emph{Hidden Markov/State-space drift models.} They capture non-stationarity but at the cost of additional latent states, transition matrices, and identifiability issues in short sessions; interviews suggested smooth drifts and occasional jitters rather than discrete regime switches.
\item \emph{Nonparametric (GPs, copulas) or black-box simulators.} These can match distributions but reduce transparency, making it harder to attribute failures to specific phenomena (contrary to our need for method stress-testing and ablation).
\item \emph{Graphical/Ising multilabel models.} They encode inter-category dependence but complicate the inclusion of institutional-group and time-varying annotator effects; the added complexity was not justified by interview-reported patterns.
\end{itemize}

\paragraph{Balance of realism, interpretability, and computational tractability.}
The chosen \emph{logistic, additively separable, hierarchical} specification directly mirrors clinicians’ narratives (thresholds on odds, additive nudges), isolates phenomena of interest for ablation (text, institutional-group effects, personal bias, drift, noise), and remains straightforward to compute. It generates binary multi-label annotations with realistic marginal prevalences, cross-category text couplings, institution-level calibration, and time-varying annotator behavior---precisely the ingredients identified during interviews---while keeping parameters interpretable and empirically estimable in subsequent sections.


\subsection{Model Formalisation}
\label{sec:model}


To evaluate our methodology under realistic
conditions, we developed a generative model of human annotation
behavior. The model captures the main sources of
variability and bias observed in real-world multi-annotator
classification tasks, while remaining tractable and interpretable.

Formally, for each text $i \in \{1,\dots,n_d\}$, category
$k \in \{1,\dots,n_c\}$ and annotator $a \in \{1,\dots,m\}$, we assume
a binary annotation $X_{i,k,a} \in \{0,1\}$ arises from a logit–additive
model:
\[
P(X_{i,k,a}=1) \;=\; \sigma\!\left(\frac{\eta_{i,k,a}}{s_{a,k}}\right),
\qquad
\sigma(x) = \frac{1}{1+e^{-x}},
\]
where $\eta_{i,k,a}$ is a latent linear predictor, and $s_{a,k}$
is an annotator--category scale factor (default $s_{a,k}=1$ in the deterministic
setting). The linear predictor is decomposed as
\[
\eta_{i,k,a} \;=\; 
\underbrace{\mu_k}_{\text{base prevalence}} +
\underbrace{S_{i,k}}_{\text{text effect}} +
\underbrace{E_{g(a),k}}_{\text{institutional-group effect}} +
\underbrace{A_{i,k,a}}_{\text{annotator $\times$ category drift}} +
\underbrace{B_{i,a}}_{\text{annotator global drift}} .
\]

\paragraph{Base prevalence.}  
The intercept $\mu_k = logit(p_{\text{base}})$ represents the
baseline tendency for category $k$ to be assigned, thus controlling
global prevalence across categories.

\paragraph{Text effect.}  
Each document has a latent representation $z_i \sim \mathcal{N}(0,I_L)$,
which interacts with category-specific loadings $\beta_{:,k} \sim
\mathcal{N}(0,\beta_{\text{scale}}^2 I)$ and a residual term
$\eta_{i,k} \sim \mathcal{N}(0,\sigma_\eta^2)$, yielding
$S_{i,k} = z_i^\top \beta_{:,k} + \eta_{i,k}$. This captures the fact
that some texts are intrinsically more likely to trigger certain
categories.

\paragraph{Institutional-group effect.}  
Annotators are grouped into institutions/services, each with
systematic biases $E_{r,k}$ on category $k$, parameterized by a
contrast $\pm \tfrac{1}{2}\,\text{school\_contrast}$. This reflects
realistic institution-level calibration differences.

\paragraph{Annotator $\times$ category drift.}  
Each annotator $a$ has idiosyncratic biases per category:
\[
A_{i,k,a} = \alpha_{0,ak} + \phi_{ak}\,u_{i,a} + \mathrm{RW}^A_{i,ak}.
\]
Here $\alpha_{0,ak} \sim \mathcal{N}(0,\alpha0\_sd^2)$ is a
fixed bias, $\phi_{ak} \sim \mathcal{N}(\phi_{\text{mean}},\phi_{\text{sd}}^2)$
models learning or fatigue as a linear function of progress
$u_{i,a}\in[0,1]$, and $\mathrm{RW}^A_{i,ak}$ is a session-reset
random walk with step size $rwA\_sd$. This construction mimics the
well-known within-annotator drifts across time.

\paragraph{Annotator global drift.}  
Similarly, annotators have global temporal tendencies across all
categories:
\[
B_{i,a} = \beta_{0,a} + \psi_a\,u_{i,a} + \mathrm{RW}^B_{i,a},
\]
with $\beta_{0,a}\sim\mathcal{N}(0,BETA0\_sd^2)$ and
$\psi_a \sim \mathcal{N}(\psi_{\text{mean}},\psi_{\text{sd}}^2)$,
together with a random walk $\mathrm{RW}^B_{i,a}$ of step size $rwB\_sd$.
This component accounts for annotators being globally stricter or
more lenient over time.

\paragraph{Hyperparameters.}  
The full parameter set governing the generative process is
\[
\Theta = \Big\{ p_{\text{base}}, \beta_{\text{scale}}, \sigma_\eta,
\;\text{school\_contrast}, \alpha0\_sd, \phi_{\text{mean}}, \phi_{\text{sd}},
\psi_{\text{mean}}, \psi_{\text{sd}}, rwA\_sd, rwB\_sd \Big\}.
\]

\medskip
This hierarchical, logit-additive model provides a principled framework
for human annotations: it reproduces the interplay of (i) intrinsic
text difficulty, (ii) systematic institutional biases, (iii) individual
annotator biases and drifts, and (iv) random temporal variability.
By tuning $\Theta$ on empirical human data, the simulator generates
synthetic annotation matrices that match observed prevalence,
agreement, and drift statistics, thus serving as a faithful yet
controllable approximation of real annotation processes.
%--------------------------------------------------------------------

\subsection{Empirical Estimation of the Parameters Range}
\label{sec:empirical-estimation}
To ensure that our simulation model does not rely on arbitrary priors, we empirically 
estimated the plausible ranges of hyperparameters from real human classifications. 
Specifically, six intensive care unit (ICU) practitioners independently applied eleven 
inclusion criteria to a set of fifty anonymized medical documents. This corpus provided a 
realistic benchmark of inter-annotator variability in a high-stakes domain.

The resulting human classification matrices (\(n=50\) documents, \(k=11\) categories, 
\(m=6\) annotators) were aggregated to compute summary statistics such as mean prevalence, 
average number of labels per text, variance across category prevalences, and pairwise 
micro-F1 between annotators. These empirical descriptors were then used as targets in an 
inverse optimization procedure: a projected gradient-descent estimation of the logistic 
model parameters (Section~\ref{sec:model}). By iteratively adjusting parameters until 
simulated classifications reproduced the empirical statistics, we obtained realistic 
estimates of the base prevalence (\(p_{\text{base}}\)), the scales of text, institutional-group, and 
annotator effects, as well as time-drift terms (\(\phi, \psi\)) and random walk 
variabilities. Bootstrap confidence intervals (percentile method) were computed across multiple random restarts to 
capture the robustness of the fit.

This procedure ensures that the hyperparameter ranges explored in our synthetic experiments 
are anchored in actual human decision patterns rather than arbitrary choices. As a result, 
the simulated distributions closely mirror the empirical variability observed among experts, 
providing a reliable basis for generating synthetic classifications in subsequent analyses.

\paragraph{Estimated ranges.}  
Across 100 restarts of the optimization procedure, we obtained the following 95\% bootstrap confidence 
intervals for the key hyperparameters (mean ± 95\% bootstrap CI):

\begin{itemize}
  \item Base prevalence: \(p_{\text{base}} \approx 0.57 \,[0.50;0.64]\).
  \item Text–category coupling: \(\beta_{\text{scale}} \approx 1.22 \,[1.17;1.27]\).
  \item Residual text noise: \(\sigma_{\eta} \approx 0.51 \,[0.49;0.53]\).
  \item Institutional-group contrast: \(\text{school\_contrast} \approx 1.08 \,[0.86;1.30]\).
  \item Annotator bias spread: \(\alpha_{0,\text{sd}} \approx 0.25 \,[0.20;0.29]\).
  \item Category-specific fatigue/learning slope spread: \(\phi_{\text{sd}} \approx 0.30 \,[0.25;0.35]\).
  \item Global drift spread: \(\psi_{\text{sd}} \approx 0.26 \,[0.21;0.31]\).
  \item Random-walk variability (category/global): \(rwA_{\text{sd}} \approx 0.045 \,[0.036;0.055]\),
        \(rwB_{\text{sd}} \approx 0.055 \,[0.046;0.065]\).
  \item Time-drift means: \(\phi_{\text{mean}} \approx -0.43 \,[-0.50;-0.36]\),
        \(\psi_{\text{mean}} \approx -0.43 \,[-0.50;-0.36]\).
\end{itemize}

These ranges confirm that expert variability is shaped not only by a baseline prevalence 
level but also by substantial annotator- and institutional-group-specific effects, together with temporal 
drifts. Notably, the negative values of \(\phi_{\text{mean}}\) and \(\psi_{\text{mean}}\) 
indicate a consistent tendency toward stricter classifications over the course of annotation 
sessions, counterbalancing the relatively high base prevalence. This interplay highlights 
the importance of modeling dynamic annotator behavior rather than relying solely on static 
bias terms.


\subsection{Simulation Computation}
\label{sec:simulation-computation}

The simulation was implemented in Python, building on the generative logistic model as described. For each combination of hyperparameters in the
exploration grid (n=177147), the program generated classification matrices for:
(i) six human annotators, (ii) six unbiased AI models (\texttt{AI\_good}), and
(iii) six biased AI models (\texttt{AI\_bad}).

\paragraph{Human and AI\_good generation.}
The human and \texttt{AI\_good} classifications were generated using the same
distributional assumptions and hyperparameters, ensuring that the AI\_good agents
differed from humans only through random draws. They therefore reflect plausible,
unbiased automated annotators operating within the human-variability space.

\paragraph{AI\_bad generation.}
In contrast, \texttt{AI\_bad} agents were deliberately biased by assigning them
distinct \emph{institutional-group effects}. In practice, each biased AI was associated with a
unique group vector $E_{r,k}$ applied across categories. The values of these
group effects were drawn with larger contrasts than those observed among humans,
and then normalized to ensure that each AI\_bad was clearly separated from the
human annotators in the latent logit space. In practice, this amounts to shifting
the log-odds for each category by an additional group-specific term of magnitude
$\Delta \in [1.0,\,2.0]$. This procedure ensures that
each \texttt{AI\_bad} systematically diverges from human-comparable behavior while still
producing coherent binary classifications.

\paragraph{Grid exploration.}
All hyperparameters were explored on a grid defined by the empirical ranges
estimated in Section~\ref{sec:empirical-estimation}, using three evenly spaced
points per parameter (lower bound, empirical mean, and upper bound). This yielded a
combinatorial set of parameter configurations, each producing a full corpus of
human and AI annotations. The entire set of results was subsequently aggregated
for performance evaluation of the discrimination algorithms (Section~\ref{sec:algorithms-eval}).

\subsection{Code and Reproducibility}
\label{sec:code-reproducibility}

To support reproducibility, we organized the computational workflow as a version-controlled public repository containing: (i) the reusable Python implementation of \textit{Rank from References}, (ii) the scripts used to generate the evaluative simulation grid and benchmark RFR against the F1-based baselines, (iii) the sensitivity-analysis workflow for panel size and quantile threshold, and (iv) the \LaTeX{} manuscript source itself \cite{rfrgithub}. In practical terms, reproducibility is structured in three layers: a package layer for the core scoring functions, a script layer for end-to-end simulation and benchmarking, and a manuscript layer for the text and figures reported in the preprint.

The repository was designed so that analyses can be rerun from the project root with explicit commands, fixed default seeds, and generated outputs written to dedicated non-versioned result directories. This separation between source code, manuscript source, and derived artifacts provides an auditable framework for reproducing the main in silico benchmark as well as the supplementary sensitivity experiments.

\subsection{Algorithm Performance Evaluation Methodology}
\label{sec:algorithms-eval}

To assess whether each discrimination algorithm is able to distinguish unbiased AI models
(\texttt{AI\_good}) from biased ones (\texttt{AI\_bad}), we conducted a systematic evaluation
over the entire hyperparameter grid generated by our simulator.
For each parameter configuration, the simulator outputs six human classifications and
twelve AI classifications (six unbiased and six biased), providing a controlled benchmark
where ground truth is explicitly known.

\paragraph{Evaluation principle.}
Each discrimination algorithm outputs a binary decision indicating whether a given AI
classification is ``unbiased'' (i.e., comparable to humans) or ``biased'' (i.e., out of the
acceptable human variability space). This prediction is compared against the known status
(\texttt{AI\_good} as positives, \texttt{AI\_bad} as negatives), yielding a correctness matrix
per configuration. Aggregating these decisions across all configurations allows computing
the classical information retrieval measures of \emph{precision}, \emph{recall}, and
\emph{F1-score} for each algorithm, with the target task being the correct identification of
\texttt{AI\_good} models. The evaluation is repeated across all hyperparameter combinations,
ensuring that performance estimates are not driven by a particular set of simulation conditions but
reflect the robustness of each method.

\paragraph{Algorithms under study.}
We implemented and compared three complementary discrimination strategies:

\begin{enumerate}
  \item \textbf{Rank from References (RFR).} 

  \item \textbf{Majority-F1.}  
  A strict majority vote among the six human annotators is used as a surrogate
  consensus. The F1-score of each actor (human or AI) against this consensus is
  computed. An AI is predicted unbiased if its F1-score is at least as high as the
  worst human annotator’s F1-score.

  \item \textbf{Average pairwise F1.}  
  Each human is iteratively set as the reference; for each actor, the F1-score against
  this reference is computed. Scores are then averaged across all references. An AI
  is predicted unbiased if its mean F1-score is greater than or equal to the worst human
  mean F1-score. This approach avoids privileging a single consensus and explicitly
  incorporates annotator plurality.

\end{enumerate}

\paragraph{Outputs.}
For each algorithm, we saved: (i) the boolean correctness matrix across all
configurations, (ii) per-configuration counts of true/false positives and negatives, and
(iii) aggregate performance metrics (precision, recall, F1-score). This comprehensive
setup enables a robust comparison of the algorithms’ ability to generalize across
different annotation regimes while remaining faithful to the diversity observed in
human experts.

\paragraph{Sensitivity protocol for panel size and quantile threshold.}
To evaluate robustness of the decision rule, we performed an additional sweep varying panel size
\(m\in\{3,4,5,6,8,10,12\}\) and non-inferiority quantile \(q\in\{0.75,0.9,1.0\}\), where \(q=1.0\)
corresponds to the most permissive quantile anchor. For each \((m,q)\) cell, 60 independent replicates were
simulated with the same data-generating framework, and uncertainty was propagated with
document-level bootstrap (\(B=200\), \(\alpha=0.05\)). We report acceptance rate of biased AIs,
specificity, recall, and threshold statistics per cell.

\subsection{In Silico Results}

Across the full hyperparameter grid ($N=2{,}125{,}764$ AI classifications), \textit{Rank from References} exhibited a more conservative operating profile than F1-based criteria. RFR prioritized specificity---i.e., rejecting biased ($AI_{bad}$) classifiers that fall outside the human-variability space---whereas F1-based rules maximized recall---i.e., accepting more unbiased ($AI_{good}$) classifiers. RFR reached \textbf{specificity $45.9\%$}, \textbf{precision $58.2\%$}, \textbf{recall/sensitivity $60.7\%$}, and \textbf{F1 $54.9\%$}; by comparison, the Majority-F1 rule achieved $20.5\%$ specificity, $53.7\%$ precision, $81.8\%$ recall, and $61.3\%$ F1, while the Average pairwise F1 rule achieved $21.9\%$ specificity, $55.3\%$ precision, $85.0\%$ recall, and $63.4\%$ F1. Overall accuracy remained comparable across methods (RFR $53.3\%$, Average pairwise F1 $53.4\%$, Majority-F1 $51.2\%$), and Matthews correlation coefficients were modest (RFR $0.066$, Average pairwise F1 $0.088$, Majority-F1 $0.030$), underscoring the difficulty of the discrimination task under broad, realistic annotator variability.

From a safety perspective, the error profiles are instructive. Relative to Average pairwise F1, RFR \emph{reduced false acceptances of biased AIs} ($FP$) by $\mathbf{254{,}810}$ (from $830{,}297$ to $575{,}487$), i.e.\ a relative reduction of about \(30.7\%\), but also \emph{accepted fewer unbiased AIs} ($TP$) by $\mathbf{257{,}937}$ (from $903{,}032$ to $645{,}095$), i.e.\ about \(28.6\%\) fewer true positives. This deliberate trade-off toward conservatism aligns with our $W$-set acceptance principle and the two-CI non-inferiority rule. In other words, RFR is more reluctant to certify an algorithm as ``human-comparable'' unless its disagreement patterns remain consistently within expert--expert divergences observed in the reference panel, thereby lowering the risk of endorsing systematically biased models.

\paragraph{Consensus percentile reporting.}
Beyond binary acceptance, we report the normalized micro-RFR score as a consensus percentile. If $\tilde{\rho}_{\chi}(x)=p$, the system can be described as behaving around the $100p$-th percentile of expert-referenced disagreement ranks. This provides an interpretable continuous layer on top of the go/no-go rule and facilitates communication with non-statistical stakeholders.

\paragraph{Sensitivity to panel size $m$ and non-inferiority quantile $q$.}
Using the protocol described above, results show that permissiveness is strongly driven by $q$: the mean acceptance rate of biased AIs ranged from $27.5\%$ to $84.4\%$ for $q=0.75$, from $43.3\%$ to $90.0\%$ for $q=0.9$, and from $62.8\%$ to $91.1\%$ for $q=1.0$. For fixed $q$, increasing $m$ did not induce a monotonic inflation in all regimes, but generally improved specificity at conservative quantiles (e.g., at $q=0.75$, specificity increased from $15.6\%$ at $m=3$ to $72.5\%$ at $m=12$, while recall decreased from $98.1\%$ to $80.3\%$). Overall, these analyses support treating $q$ as the primary clinical lever and interpreting $m$ effects through explicit sensitivity checks rather than assumptions of fixed permissiveness.

\paragraph{Performance variation across model configurations.}
Across the explored hyperparameter grid, performance metrics showed only small-amplitude fluctuations. Although several pairwise comparisons reached statistical significance---unsurprising given the large number of simulated decisions---the corresponding effect sizes were not clinically meaningful. The largest primary-metric deviation we observed was a modest F1 decrease of 0.05 for \textit{RFR} when increasing institutional-group contrast (\texttt{school\_contrast}) from 0.860 to 1.300; all other shifts remained within smaller narrow bands. Taken together, these results indicate that the in silico conclusions are stable with respect to reasonable hyperparameter perturbations, with no clinically relevant degradation in performance across model configurations.

\paragraph{Practical takeaway.} For high-throughput screening, Average pairwise F1 offers a higher yield (recall) of candidate unbiased AIs. However, for \emph{go/no-go} decisions in safety-critical settings, RFR’s higher specificity and precision make it a more suitable gatekeeper: it preferentially rejects out-of-space behaviors and accepts only those classifiers whose disagreement envelopes remain \emph{within} credible expert variability.






\section{Discussion}

\subsection{Positioning \textit{Rank from References} within the state of the art and scope for generalization}

\paragraph{Relation to existing families.}
Work on disagreement and imperfect supervision broadly spans three strands. First, \emph{latent-label inference} and annotator modeling estimate unobserved “true” labels from discordant raters (e.g., Dawid–Skene and its diagnostic-test analogs) \cite{dawid1979maximum,hui1980study}, as well as Bayesian formulations that place priors on annotator reliability or smooth the empirical label distribution to reflect uncertainty \cite{chou2019,fornaciari2021}. Second, \emph{learning with noisy labels} focuses on training-time robustness (co-teaching, bootstrapping, consistency regularization), synthesizing theoretical and empirical advances \cite{garg2021,han2022,han2023,frenay2014,cannings2018}. Third, \emph{embedding and nearest-neighbor} approaches evaluate or retrieve by distances in representation spaces rather than by exact matches \cite{disalvo2024}. While these families are valuable, they seldom offer an operational, auditable \emph{go/no-go} rule for deployment when truth is \emph{plural} and disagreement among competent experts is expected—a situation extensively discussed in perspectivist NLP \cite{basile2021} and in patient-safety research that cautions against brittle single-number conclusions \cite{shojania2020}.

\textit{Rank from References} (RFR) addresses this gap by (i) constructing an empirically grounded human-variability envelope $W$ from expert--expert distances, (ii) ranking a candidate system against each reference via a simple distance statistic, and (iii) deciding acceptance through a non-inferiority ladder with confidence intervals, eschewing extra thresholds. In spirit, this is closer to a nonparametric comparison of discrepancy ranks than to consensus relabeling, akin to rank-based reasoning used for discontinuous or discrete outcomes \cite{conover1972}. Unlike standard F1 checks against an inferred consensus, RFR preserves the \emph{plurality} of competent judgments and evaluates whether the system operates \emph{within} that envelope.

\paragraph{Why RFR fits our scope.}
Our application is multi-label, document-level, binary classification with non-exclusive clinical categories and multiple credible annotators. In this regime, distances over $\{0,1\}^{n_d\times n_c}$ (e.g., Hamming) are transparent, easy to audit, and map directly to “bits flipped” across documents and categories—properties that clinicians readily interpret. RFR’s acceptance rule, anchored in human–human variability, complements noisy-label training remedies \cite{han2022,han2023,frenay2014} and avoids collapsing disagreement into a single “gold” label \cite{dawid1979maximum,hui1980study,chou2019,fornaciari2021}. The emphasis on specificity to biased behavior aligns with safety-first evaluation in clinical AI \cite{shojania2020} and with recent perspectives advocating logically consistent, variance-aware assessment of LLM outputs \cite{SLCA}.

\paragraph{Toward generalized RFR for rich text structurations.}
Clinical text extraction frequently targets \emph{structured} outputs (e.g., FHIR-like forms) with many semantically related fields. In these settings, exact-match metrics can be misleading by penalizing near-miss codes as harshly as semantically distant errors. The RFR principle should extend naturally: retain the rank-based acceptance logic, but generalize the base distance $d(\cdot,\cdot)$ to be \emph{structure- and semantics-aware}. Options include cost-weighted set distances or ontology-path distances for hierarchical vocabularies, graph/optimal-transport–style discrepancies for many-to-many correspondences, and fieldwise compositions for heterogeneous forms—thus granting partial credit to semantically proximal alternatives and sharper penalties to distant ones. This mirrors the move from hard labels to perspective-aware distributions \cite{basile2021}, leverages robust distance design familiar to the noisy-label literature \cite{han2022,han2023,frenay2014}, and resonates with nearest-neighbor reasoning in high-dimensional embeddings when appropriate \cite{disalvo2024}. In short, a \emph{generalized RFR} with clinically informed distances could evaluate arbitrary text structurations with large category sets, where semantic proximity matters—avoiding the pitfalls of exact matching while preserving an auditable, parameter-light acceptance decision.



\subsection{Strengths}

Our study has several strengths spanning design, methodology, and practical relevance.

\paragraph{Clinically grounded simulation.}
The simulator was constructed from semi-structured interviews with intensivists after a real labeling session and encodes the specific phenomena clinicians reported: heterogeneous base rates, text-driven co-activations, institution-level calibration, annotator-specific thresholds, and within-shift drift. This qualitative-to-quantitative pipeline yields synthetic data that preserve the qualitative texture of clinical annotation while remaining fully controllable for stress testing.

\paragraph{Interpretable, principled model structure.}
The chosen logistic, additively separable specification provides direct odds-ratio interpretability and a transparent decomposition of effects (text, institutional-group, annotator$\times$category, global annotator). Each term corresponds to an observable source of variability rather than a black-box proxy, enabling clear attribution of failure modes and targeted ablations. Importantly, all hyperparameters are used to express phenomena identified by clinicians; their values are estimated empirically rather than fixed arbitrarily, preserving interpretability without sacrificing flexibility.

\paragraph{Safety-aligned evaluation criterion.}
Rank from References (RFR) operationalizes a conservative acceptance rule that privileges \emph{specificity} to biased behavior over marginal gains in recall. By anchoring acceptance to the envelope of human–human disagreement and validating decisions with confidence-interval–based non-inferiority, the framework aligns with safety-first principles that are essential in clinical decision support. The resulting error profile---which reduces false acceptances of biased systems---better matches the asymmetric costs in real deployment.

\paragraph{Calibration without new thresholds.}
RFR relies on inter-expert disagreement structure (e.g., Hamming distances) rather than ad hoc probability thresholds or additional hyperparameters. This reduces tuning burden and the risk of overfitting acceptance criteria to a particular dataset or prevalence spectrum, facilitating reproducible evaluation across sites and time.

\paragraph{Robustness to reasonable perturbations.}
In extensive in silico experiments, performance differences across hyperparameter configurations were small in magnitude and not clinically consequential. This stability supports external validity of the conclusions and suggests that the comparative behavior of methods (e.g., RFR vs.\ F1-based rules) is preserved across plausible shifts in institutional calibration or annotator drift intensity.

\paragraph{Faithful multi-label, temporal dynamics.}
The simulator natively handles multi-label outputs, inter-category coupling via latent text factors, and both category-specific and global annotator drifts (linear trends and session-reset random walks). This combination is uncommon in existing annotation models and is crucial for tasks where categories co-activate and annotator behavior evolves within a shift.

\paragraph{Ablation-ready and diagnosis-friendly.}
Because the components of the linear predictor map to distinct real-world phenomena, the setup supports clean ablations (e.g., removing institutional-group effects, suppressing drift, or turning off text loadings) to isolate which sources of variability stress a method. This design shortens the path from ``observed failure'' to ``actionable hypothesis'' about its cause.

\paragraph{Generalizable and site-portable.}
Institution effects and annotator-specific terms make domain shift an explicit part of the data-generating process, rather than an afterthought. Methods that perform well here are more likely to maintain behavior under changes in staffing, local practice patterns, or documentation styles, supporting portability across hospital services.

\paragraph{Operational readiness and auditability.}
The evaluation pipeline produces decisions that can be traced back to human disagreement envelopes and CI-based rules, simplifying communication with clinical partners and governance bodies. Decisions are reproducible, explainable, and auditable—practical prerequisites for adoption in safety-critical environments.

\paragraph{Easy computability at scale.}
Despite accounting for a variety of concurrent processes, the framework remains lightweight: binary outcomes, a logistic link, and additive components allow efficient simulation over large grids and rapid recomputation for sensitivity analyses. This easy computability enables broad exploration of settings and supports routine regression testing as systems evolve.

\paragraph{Method-agnostic benchmarking.}
Finally, the benchmark decouples data generation (human-like variability) from acceptance logic (RFR vs.\ baselines), allowing fair, apples-to-apples comparisons among alternative evaluation strategies. By avoiding reliance on any single performance metric or threshold, the framework encourages more nuanced, safety-aware method development.

\subsection{Limitations}

\paragraph{Scope and elicitation base.}
The simulator is qualitatively grounded in interviews following a single pilot session ($n=50$ documents, $k=11$ categories, $m=6$ annotators). While this provides authentic clinical texture, it constrains external validity: the elicited phenomena (e.g., drift intensity, institutional calibration) may not fully represent other services, pathologies, documentation styles, or multilingual settings.

\paragraph{Modeling simplifications.}
We adopt a logistic, additively separable linear predictor with Gaussian latent factors for text and simple contrasts for institutions. This choice emphasizes interpretability but omits richer structures: non-additive interactions among components, non-Gaussian or heavy-tailed text effects, and explicit inter-category graphical dependencies (e.g., Ising/conditional random fields). As such, some higher-order co-activation patterns or rare decision heuristics may be underrepresented.

\paragraph{Drift representation.}
Annotator dynamics are captured via linear trends and session-reset random walks. These capture smooth evolution and short-term jitter but do not model discrete regime switches, circadian effects, break-related resets, or workload-dependent thresholds. Consequently, abrupt behavioral shifts or complex temporal autocorrelation structures could be mis-specified.

\paragraph{Binary outcome formulation.}
Categories are simulated as binary decisions. Although this mirrors the annotation interface, several clinical constructs are inherently graded or ordinal. The binary formulation may compress clinically meaningful severity distinctions and attenuate calibration challenges tied to borderline cases.

\paragraph{Identifiability and parameter estimation.}
Multiple variance components (text, institutional-group, annotator-by-category, global annotator) can trade off in finite samples, raising partial identifiability concerns. Hyperparameters are estimated empirically but, given the modest elicitation base, may absorb dataset-specific artifacts (e.g., documentation templates), limiting portability.

\paragraph{Institution effect simplification.}
The institutional-group effect is parameterized via a simple contrast per category. Real institutional differences can be multi-faceted (coding policy, staffing mix, supervision practices). A single-contrast summary may underfit such heterogeneity, potentially smoothing over site-specific quirks that matter in practice.

\paragraph{Inter-category dependence.}
Co-activation is induced through shared text factors, not an explicit category–category dependency structure. This may understate conditional dependencies that persist after text effects are accounted for (e.g., mutually exclusive labels or strict hierarchical relations).

\paragraph{Metric and objective emphasis.}
While we report a broad set of metrics, the study emphasizes specificity-oriented acceptance (RFR) versus recall-oriented baselines. This lens matches safety-first goals but may understate scenarios where high sensitivity to rare but benign behaviors is preferred (e.g., exploratory triage tools). Moreover, clinically “non-relevant” effect sizes at the aggregate level could still matter in subpopulations or edge cases not fully explored.

\paragraph{Dependence on reference panel quality.}
RFR calibrates acceptance to the envelope of human–human disagreement, which is the foundational premise of the framework. Operationally, if the reference panel encodes systematic biases or inconsistent practices, the acceptance boundary may inherit these features—potentially rejecting models that challenge suboptimal norms or accepting models aligned with them.

\paragraph{Distance choice and label topology.}
Our disagreement geometry relies on Hamming distances across categories. Alternative notions of proximity (e.g., cost-sensitive, hierarchy-aware, or clinically weighted distances) might better reflect the asymmetries among labels; using Hamming may misstate the severity of certain disagreements.

\paragraph{Simulation–reality gap.}
Despite qualitative alignment, any simulator is an approximation. Documentation shifts over time, EHR interface changes, guideline updates, and staffing turnover can all induce distribution drift not fully captured by the current parameterization. Robustness observed in silico does not guarantee invariance to such real-world evolution.

\paragraph{Computational and stochastic considerations.}
Large grids enable sensitivity analysis but remain subject to Monte Carlo variability and seed dependence; while we observe narrow performance bands, extremely rare configurations or adversarial settings may be under-sampled. Additionally, we do not propagate full Bayesian uncertainty over hyperparameters into downstream acceptance summaries.

\paragraph{Privacy and data provenance.}
The qualitative foundation relies on clinician recall and behavior within a specific institutional context. Although no patient-identifying data are simulated or reported, nuances stemming from local documentation habits may implicitly influence the learned hyperparameters, raising transferability questions when moving to hospitals with different governance or note-taking cultures.




\subsection{Conclusion}

Situated alongside latent-label aggregation and annotator modeling \cite{dawid1979maximum,hui1980study,chou2019,fornaciari2021}, robustness-to-noise training \cite{garg2021,han2022,han2023,frenay2014,cannings2018}, and representation-based comparators \cite{disalvo2024}, \textit{Rank from References} offers a deployment-oriented alternative: it asks whether a system’s disagreements remain \emph{within} empirically observed human variability, providing an auditable, parameter-light go/no-go rule rather than a single “gold” label score \cite{basile2021,shojania2020,conover1972,SLCA}. In our scope---multi-label, non-exclusive, document-level clinical text---this plurality-respecting criterion complements training-time robustness and yields decisions aligned with safety-first practice. Looking forward, the same acceptance logic extends naturally to richer text structuring targets (e.g., FHIR-like forms) by replacing Hamming with clinically informed, semantics-aware distances, thus crediting near misses and penalizing distant errors while preserving a clear acceptance boundary \cite{basile2021,disalvo2024}. Accordingly, we contend that variability-aware acceptance metrics of this kind should be integrated into pre-deployment evaluation of AI text-classification systems that can impact patient care \cite{shojania2020,SLCA}.










\section{Declarations}

\subsection{Ethics Statement}
This manuscript reports a methodological study whose benchmark results and sensitivity analyses are generated from simulated data. The empirical calibration step relied on secondary use of anonymized source documents, expert annotations, and high-level qualitative themes derived from a prior study evaluating AI engines for the classification of intensive care reports against inclusion criteria. That primary study was approved by the Institutional Review Board (IRB) and the Scientific and Ethical Committee (CSE) of the Montpellier University Hospital Health Data Warehouse (eDOL) (IRB No.\ A034/2024-10-152/001, approved on 28 November 2024), and the secondary use reported here was covered by that protocol. No patient-identifying information is reported in the manuscript or distributed in the accompanying repository. No prospective enrollment, intervention, or clinical trial was performed for the present methodological study, and clinical trial registration was therefore not applicable. All source material was handled under applicable institutional confidentiality and data-protection requirements.

\subsection{Funding Statement}
This work received no external project-specific funding. It was conducted within the routine salaried institutional activities of the authors at CHU Montpellier and Univ. Montpellier.

\subsection{Competing Interest Statement}
The authors declare that they have no competing interests.

\subsection{Data Availability Statement}
The study is reproducible from the public code repository \cite{rfrgithub}. The results reported in the manuscript arise from a generative simulation framework and can be regenerated by rerunning the documented workflows in the repository. No patient-level source documents are redistributed. The manuscript reports only simulated outputs and aggregate calibration summaries derived from approved secondary use of anonymized source material.

\bibliographystyle{plain}
\begin{thebibliography}{9}
\bibitem{dawid1979maximum}
Dawid, A. P., and Skene, A. M. (1979). Maximum likelihood estimation of observer error-rates using the EM algorithm. \textit{Applied Statistics}, 28(1), 20--28.

\bibitem{hui1980study}
Hui, S. L., and Walter, S. D. (1980). Estimating the error rates of diagnostic tests. \textit{Biometrics}, 167--171.

\bibitem{chou2019}
Chou, H., and Lee, M. (2019). Bayesian aggregation for imperfect annotations: Leveraging prior distributions of annotator reliability. \textit{International Journal of Approximate Reasoning}, 105, 25--37.

\bibitem{fornaciari2021}
Fornaciari, T., Moschitti, A., and Jha, R. (2021). Expressing annotator disagreement with Bayesian label smoothing. \textit{Transactions of the Association for Computational Linguistics}, 9, 1234--1250.

\bibitem{garg2021}
Garg, S., et al. (2021). \textit{Robust learning from noisy labels with co-teaching and dynamic bootstrapping}. In IEEE Conference on Computer Vision.

\bibitem{disalvo2024}
Di Salvo, V., et al. (2024). WANN: Weighted approximate nearest neighbors for classification in large-scale embeddings. \textit{Journal of Machine Learning Research}, 25, 1--29.

\bibitem{basile2021}
Basile, V., et al. (2021). Toward a perspectivist data manifesto. \textit{Computational Linguistics}, 47(3), 649--674.

\bibitem{han2022}
Han, B., et al. (2022). Learning with imperfect data at scale: challenges and solutions. \textit{IEEE Transactions on Pattern Analysis and Machine Intelligence}, 44(5), 2500--2513.

\bibitem{han2023}
Han, B., et al. (2023). Imperfect data learning: a comprehensive overview. \textit{arXiv preprint} arXiv:2302.12345.

\bibitem{frenay2014}
Fr\'enay, B., and Verleysen, M. (2014). Classification in the presence of label noise: a survey. \textit{IEEE Transactions on Neural Networks and Learning Systems}, 25(5), 845--869.

\bibitem{cannings2018}
Cannings, T., Nath, A., and Sen, B. (2018). On the impact of label noise on the complexity and convergence rates of Empirical Risk Minimization. \textit{NeurIPS}, 31.

\bibitem{shojania2020}
Shojania, K., and Marang-van de Mheen, P. J. (2020). Challenges in making reliable conclusions in patient safety research. \textit{BMJ Quality \& Safety}, 29(4), 351--354.

\bibitem{conover1972}
Conover, W.\ J.\ (1972). A Kolmogorov Goodness-of-Fit Test for Discontinuous Distributions. \textit{Journal of the American Statistical Association}, 67(339), 591--596.

\bibitem{SLCA}
Loi et al. (2025). Self-Logical Consistency Assessment of Large Language Models for Patient Feedback Classification: Algorithm Development and Validation Study. \textit{medRxiv}.

\bibitem{who2025lmm}
World Health Organization (2025). Ethics and governance of artificial intelligence for health: Guidance on large multi-modal models (LMMs). \textit{World Health Organization}, Geneva. 

\bibitem{fda2025draft}
U.S. Food and Drug Administration (2025). Artificial Intelligence-Enabled Device Software Functions: Lifecycle Management and Marketing Submission Recommendations (Draft Guidance). \textit{FDA Guidance Document}. 

\bibitem{nhs2024eval}
NHS England (2024). Planning and implementing real-world artificial intelligence (AI) evaluations: lessons from the AI in Health and Care Award. \textit{NHS England Publication}. 

\bibitem{futureai2025}
Lekadir, K., et al. (2025). FUTURE-AI: International consensus guideline for trustworthy and deployable artificial intelligence in healthcare. \textit{BMJ}, 388, e081554. 

\bibitem{tam2024humanEval}
Tam, T.\ Y.\ C., et al. (2024). A framework for human evaluation of large language models in healthcare. \textit{npj Digital Medicine}, 7, 258. 

\bibitem{singhal2023clinKnowledge}
Singhal, K., et al. (2023). Large language models encode clinical knowledge. \textit{Nature}, 620, 172--180. 

\bibitem{reddy2024genAI}
Reddy, S., et al. (2024). Generative AI in healthcare: opportunities, risks and implementation considerations. \textit{Implementation Science}, 19, 62.

\bibitem{rfrgithub}
Loi, Z. (2026). \textit{rank-from-references}. GitHub repository. \url{https://github.com/ZenoLoi/rank-from-references}




\end{thebibliography}




\end{document}
